\documentclass[11pt,a4paper]{article}

% --- Packages ---
\usepackage[utf8]{inputenc}
\usepackage[T1]{fontenc}
\usepackage[french]{babel}
\frenchsetup{StandardLists=true}
\usepackage[margin=1.5cm,bottom=0.8cm]{geometry}
\usepackage{fontawesome5} 
\usepackage{xcolor} 
\usepackage{titlesec} 
\usepackage{enumitem} 
\usepackage{tabularx} 
\usepackage[hidelinks]{hyperref} 

% --- Custom Colors ---
\definecolor{graytext}{HTML}{7A7A7A}

% --- Section Formatting ---
\titleformat{\section}
  {\Large\bfseries}
  {}
  {0em}
  {\rule{2cm}{3pt}\hspace{0.5em}}
\titlespacing*{\section}{0pt}{1ex}{1ex}

% --- Global Parameters ---
\setlength{\parindent}{0pt}
\pagestyle{empty}
\setlist[itemize]{label=\textbullet, leftmargin=0.4cm, topsep=0pt, parsep=0pt, partopsep=0pt}

\begin{document}

% ================= HEADER =================
\begin{minipage}[b]{0.6\textwidth}
    {\Huge \textbf{Gabriel Dahan}}
\end{minipage}
\hfill
\begin{minipage}[b]{0.35\textwidth}
    \raggedleft \small
    Rennes, FRANCE \\
    \faPhone \ +33(0) 7 61 74 04 92 \\
    \faEnvelope \ \href{mailto:work@gabrieldahan.me}{work@gabrieldahan.me} \\
    \faLinkedin \ \href{https://www.linkedin.com/in/gabriel-dahan/}{gabriel-dahan} \\
    \faGithub \ \href{https://github.com/gabriel-dahan}{gabriel-dahan} \\
    \faGlobe \ \href{https://gabrieldahan.me/}{gabrieldahan.me}
\end{minipage}

\vspace{0.4cm}

% --- Sous-titre ---
\textbf{Etudiant ingénieur à l'ENSAI, passionné par le développement et l'analyse de données}

\vspace{0.25cm}

% ================= EDUCATION =================
\section*{Formation}

\begin{tabularx}{\textwidth}{@{}>{\raggedright\arraybackslash}p{3.15cm} X@{}}
2025--12/2028 & \textbf{Diplôme d'ingénieur}, ENSAI, Rennes
  \begin{itemize}[leftmargin=0.4cm, nosep, topsep=1pt]
    \item Accrédité par la Commission des Titres d'Ingénieur (CTI). Diplôme d'ingénieur en statistiques et sciences des données (prévu).
  \end{itemize} \\[0.45em]

2025--2026\newline (1 an) & \textbf{Licence de Mathématiques}, Université de Rennes I, Rennes
  \begin{itemize}[leftmargin=0.4cm, nosep, topsep=1pt]
    \item \textbf{Spécialisation :} Mathématiques et interactions, avec obtention de 180 crédits ECTS à l'issue du cursus.
  \end{itemize} \\[0.45em]

2023--2025\newline (2 ans) & \textbf{Classes préparatoires aux grandes écoles (CPGE)}, Lycée Charlemagne, Paris
  \begin{itemize}[leftmargin=0.4cm, nosep, topsep=1pt]
    \item Préparation intensive en mathématiques avancées et informatique, avec acquisition de la gestion du temps, de la résilience et de la discipline sous une forte charge de travail.
  \end{itemize}
\end{tabularx}

% ================= EXPERIENCE =================
\section*{Expérience}

\begin{tabularx}{\textwidth}{@{}>{\raggedright\arraybackslash}p{3.15cm} X@{}}
10/2026 & \textbf{Hackathon : Build Challenge 2}, Apple × MWM × Unaite × Cmd+F, Paris
  \begin{itemize}[leftmargin=0.4cm, itemsep=1pt, topsep=0pt]
    \item Développement de \textbf{Jefe}, application de cuisine gamifiée (défis photo quotidiens, streaks et XP), en compétition parmi 10 équipes.
  \end{itemize} \\

05/2026--08/2026\newline (3 mois) & \textbf{Stage : Développement de produits analytiques}, LiveRamp, Paris
  \begin{itemize}[leftmargin=0.4cm, itemsep=1pt, topsep=0pt]
    \item Conception de fonctionnalités analytiques sur Google Cloud Platform pour valoriser les données à grande échelle.
  \end{itemize} \\
\end{tabularx}

% ================= ACADEMIC PROJECTS =================
\section*{Projets académiques}
\begin{itemize}[leftmargin=0.4cm, itemsep=2pt]
    \item \textbf{Projet de statistiques} -- Réalisation d'une étude sur le thème « Technologie et performances éducatives ».
    \textbf{Technologies :} R, \LaTeX, Quarto.
    \item \textbf{Framework de données extensible (ORM)} -- Sportix : Conception d'un ORM léger s'adaptant à tout jeu de données dès lors que sa classe de mapping relationnel est définie.
    \textbf{Technologies :} Python, pandas.
    \item \textbf{Projet d'économie} -- Réalisation d'une étude sur le thème « Obsolescence des compétences et emploi ».
    \textbf{Technologies :} Quarto, \LaTeX.
    \item \textbf{TIPE : Génération de terrains} -- Recherche et implémentation d'algorithmes (bruit de Perlin et flou par filtre moyen) pour générer des cartes de hauteur réalistes, avec visualisation de l'impact des paramètres, optimisation des algorithmes et analyse de complexité.
    \textbf{Technologies :} Python, OCaml.
\end{itemize}

% ================= PERSONAL PROJECTS =================
\section*{Projets}
\begin{itemize}[leftmargin=0.4cm, itemsep=2pt]
    \item \textbf{\href{https://comx.click/}{COMX}} -- Développement d'une plateforme de lecture de comics numériques alliant monétisation, API ouverte et traduction automatique pour une diffusion multilingue.
    \textbf{Technologies :} VueJS, Remult, Pinia, TailwindCSS, DaisyUI, TypeScript.
    \item \textbf{VaderMap} -- Développement d'un site présentant l'art d'Invader, avec enregistrement de la position des mosaïques trouvées et de leur état (intactes, cassées, etc.).
    \textbf{Technologies :} Python, Flask, JavaScript, SQLite. 
\end{itemize}

% ================= SKILLS =================
\section*{Compétences}
\begin{itemize}[leftmargin=0.4cm, itemsep=1pt]
    \item \textbf{IT \& Programmation :} Python (NumPy, pandas, scikit-learn, Matplotlib, Flask, SQLAlchemy, psycopg2), Javascript (NodeJS, Express, Typescript, VueJS), API REST, JSON, HTML/CSS, OCaml, Git.
    \item \textbf{Data, Cloud \& SGBD :} R, SQL, Google Cloud Platform (GCP), Cloudflare R2, Metabase, BigQuery, PostgreSQL, SQLite.
    \item \textbf{Édition scientifique :} \LaTeX, Markdown, Quarto, Jupyter.
    \item \textbf{Langues :} Français (Natif), Anglais (Avancé), Espagnol (Intermédiaire), Japonais (Débutant).
\end{itemize}

\end{document}